% !TEX TS-program = pdflatex
% !TEX encoding = UTF-8 Unicode

% This file is a template using the "beamer" package to create slides for a talk or presentation
% - Giving a talk on some subject.
% - The talk is between 15min and 45min long.
% - Style is ornate.

% MODIFIED by Jonathan Kew, 2008-07-06
% The header comments and encoding in this file were modified for inclusion with TeXworks.
% The content is otherwise unchanged from the original distributed with the beamer package.

\documentclass{beamer}
\setbeamercovered{invisible}



% Copyright 2004 by Till Tantau <tantau@users.sourceforge.net>.
%
% In principle, this file can be redistributed and/or modified under
% the terms of the GNU Public License, version 2.
%
% However, this file is supposed to be a template to be modified
% for your own needs. For this reason, if you use this file as a
% template and not specifically distribute it as part of a another
% package/program, I grant the extra permission to freely copy and
% modify this file as you see fit and even to delete this copyright
% notice. 


\mode<presentation>
{
  \usetheme{Warsaw}
  % or ...

  % or whatever (possibly just delete it)
}


\usepackage[english]{babel}
% or whatever

\usepackage[utf8]{inputenc}
% or whatever

\usepackage{times}
\usepackage[T1]{fontenc}
% Or whatever. Note that the encoding and the font should match. If T1
% does not look nice, try deleting the line with the fontenc.

%%% MATH RELATED

\input{macros.tex}

\title{MATH 350-2 Advanced Calculus} 
\subtitle
{} % (optional)

\author[W.R. Casper] % (optional, use only with lots of authors)
{W.R. Casper}
% - Use the \inst{?} command only if the authors have different
%   affiliation.

\institute[California State University Fullerton] % (optional, but mostly needed)
{
  Department of Mathematics\\
  California State University Fullerton}
% - Use the \inst command only if there are several affiliations.
% - Keep it simple, no one is interested in your street address.

\subject{Talks}
% This is only inserted into the PDF information catalog. Can be left
% out. 



% If you have a file called "university-logo-filename.xxx", where xxx
% is a graphic format that can be processed by latex or pdflatex,
% resp., then you can add a logo as follows:

% \pgfdeclareimage[height=0.5cm]{university-logo}{university-logo-filename}
% \logo{\pgfuseimage{university-logo}}



% Delete this, if you do not want the table of contents to pop up at
% the beginning of each subsection:
\AtBeginSubsection[]
{
  \begin{frame}<beamer>{Outline}
    \tableofcontents[currentsection,currentsubsection]
  \end{frame}
}


% If you wish to uncover everything in a step-wise fashion, uncomment
% the following command: 

%\beamerdefaultoverlayspecification{<+->}


\begin{document}

\begin{frame}
  \titlepage
\end{frame}

\begin{frame}{Outline}
  \tableofcontents
  % You might wish to add the option [pausesections]
\end{frame}

% Since this a solution template for a generic talk, very little can
% be said about how it should be structured. However, the talk length
% of between 15min and 45min and the theme suggest that you stick to
% the following rules:  

% - Exactly two or three sections (other than the summary).
% - At *most* three subsections per section.
% - Talk about 30s to 2min per frame. So there should be between about
%   15 and 30 frames, all told.

\section{Real Analysis Lecture 7}
\subsection{More on component intervals}

\begin{frame}{Component intervals}
Let $U\subseteq \mathbb{R}$ be open.
\pause
\begin{defn}
A \textbf{component interval} of $U$ is an interval $I$ with $I\subseteq U$ and with the property that if $J$ is an interval and $I\subsetneq J$, then $J\nsubseteq U$.
\end{defn}
\begin{itemize}
\pause
\item $(0,1)$ is a component interval of $\mathbb{R}\diff \{0,1,2,3\}$
\pause
\item $(-\infty,0)$ is also a component interval of $\mathbb{R}\diff \{0,1,2,3\}$
\pause
\item we will show all open sets of $\mathbb{R}$ are made of component intervals!
\end{itemize}
\end{frame}

\begin{frame}{Component intervals}
\begin{thm}[Apostol 3.10]
If $U\subseteq\mathbb{R}$ is open and $x\in U$, then there is a unique component interval of $U$ containing $x$.
\end{thm}
\pause
\begin{proof}
Uniqueness follows from previous Lemma, so we only need existence.\\
\pause
Suppose that $x\in U$ and consider
$$A = \{a: (a,x)\subseteq U\},\quad\text{and}\quad B = \{b: (x,b)\subseteq U\}$$
\pause
If $A$ is not bounded below, let $a=-\infty$.  Otherwise, let $a = \inf (A)$.\\
\pause
If $B$ is not bounded above, let $b=\infty$.  Otherwise, let $b = \sup (B)$.\\
\pause
Claim: $(a,b)$ is a component interval of $U$ containing $x$.
\end{proof}
\end{frame}

\begin{frame}{Challenge}
\begin{prob}
Prove that $(a,b)$ is a component interval of $U$.
\end{prob}
\end{frame}

\begin{frame}{Open subsets of $\mathbb{R}$}
\begin{thm}[Representation Theorem for Open Sets in $\mathbb{R}$]
If $U\subseteq\mathbb{R}$ is open, then $U$ is the union of a countable family of disjoint open intervals.
\end{thm}
\pause
\begin{proof}
From the previous theorem, each $x\in U$ is contained in a unique component interval $I_x$ of $U$.\\
\pause
Therefore
$$U = \bigcup_{x\in U} I_x$$
is a union of open intervals!\\
\pause
{\color{red}Uh oh ... this isn't a \textbf{countable} union.}
\end{proof}
\end{frame}

\begin{frame}{Open subsets of $\mathbb{R}$}
\begin{thm}[Representation Theorem for Open Sets in $\mathbb{R}$]
If $U\subseteq\mathbb{R}$ is open, then $U$ is the union of a countable family of disjoint open intervals.
\end{thm}
\begin{proof}
Instead, we consider rational points $U\cap\mathbb{Q}$.\\
\pause
If $x\in U$, then $x\in I_x\subseteq U$.\\
\pause
Any interval contains a rational number $r\in I_x$.\\
\pause
Since different component intervals must be disjoint, $I_x=I_r$.\\
\pause
Thus $x\in \bigcup_{r\in U\cap\mathbb{Q}} I_r\subseteq U$.\\
\pause
Since $x$ was arbitrary, 
$$U= \bigcup_{r\in U\cap\mathbb{Q}} I_r$$
\end{proof}
\end{frame}



\subsection{Closed Sets}

\begin{frame}{Closed sets}
\begin{defn}
A set $A\subseteq\mathbb{R}^n$ is called \vocab{closed} if it is the complement of an open set, or equivalently if its complement $\mathbb{R}^n\backslash A$ is open.
\end{defn}
\pause
\textbf{Examples:}
\begin{itemize}
\pause
\item singleton sets!
\pause
\item products of closed intervals
\end{itemize}
\end{frame}

\begin{frame}{Challenge!}
\begin{prob}
Prove that a singleton set
$$A = \{\vec a\}$$
in $\mathbb{R}^n$ is closed.
\end{prob}
\pause
\begin{soln}
We must show $U = \mathbb{R}^n\backslash\{\vec a\}$ is open.\\\pause
Suppose $\vec x\in U$.\ \pause
We must show $\vec x$ is an interior point of $U$.\\\pause
Take $r = |\vec x-\vec a|$.\ \pause Then $B(\vec x; r)$ does not contain $\vec a$.\\\pause
Therefore $B(\vec x; r)\subseteq U$, so that $\vec x$ is an interior point.\\\pause
Since $\vec x$ was arbitrary, this proves $U$ is open.
\end{soln}
\end{frame}

\begin{frame}{Challenge!}
\begin{prob}
Prove that the \textbf{closed square}
$$[a,b]\times [c,d] = \{(x,y)\in\mathbb{R}^2: a \leq x \leq b,\ c\leq y\leq d\}$$
is a closed set.
\end{prob}
\pause
\begin{soln}
$\mathbb{R}^2\backslash ([a,b]\times [c,d])$ is the same as
$$((-\infty,a)\times \mathbb{R})\cup ((b,\infty)\times \mathbb{R})\cup(\mathbb{R}\times (-\infty,c))\cup (\mathbb{R}\times (d,\infty))$$
\pause
Union of open sets is open!
\end{soln}
\end{frame}



\begin{frame}{Adherent and Accumulation Points}
\begin{defn}
A point $\vec x\in \mathbb{R}^n$ is called an \textbf{adherent point} of $A$ if for all $r > 0$ the ball $B(\vec x; r)$ contains at least one element of $A$.\\
\pause
It is called an \textbf{accumulation point} of $A$ if for all $r > 0$ the ball $B(\vec x; r)$ contains at least one element of $A$ {\color{red} different from $\vec x$}.
\end{defn}
\pause
\textbf{Examples:}
\begin{itemize}
\pause
\item every point in $A$ is adherent
\pause
\item accumulation points are adherent points
\pause
\item $-1$ is an accumulation point of $(-1,1)$.
\pause
\item suprema and infima are accumulation points!
\pause
\item $0$ is an accumulation point of $\{1/1,1/2,1/3,\dots\}$
\end{itemize}
\end{frame}

\begin{frame}{Characterizing accumulation points}
\begin{thm}[Apostol Theorem 3.17]
A point $\vec x$ is an accumulation point of $A$ if and only if for all $r > 0$, the ball $B(\vec x; r)$ contains infinitely many points of $A$.
\end{thm}
\pause
\begin{proof}
\pause
Obviously, if for all $r > 0$, the ball $B(\vec x; r)$ contains infinitely many points of $A$, then it contains at least one point of $A$ different from $\vec x$, so it's an accumulation point.\\
\pause
Suppose instead that $\vec x$ is an accumulation point and let $r > 0$.\\
\pause
It suffices to show that the set $C = (B(\vec x; r)\backslash\{\vec x\})\cap A$ is infinite.\\
\end{proof}
\end{frame}

\begin{frame}{Characterizing accumulation points}
\begin{thm}[Apostol Theorem 3.17]
A point $\vec x$ is an accumulation point of $A$ if for all $r > 0$, the ball $B(\vec x; r)$ contains infinitely many points of $A$.
\end{thm}
\pause
\begin{proof}
\pause
Suppose that $C$ is finite.\\
\pause
Then the set $\{|\vec y - \vec x|: \vec y\in C\}$ is also finite.\\
\pause
It is also nonempty and has only positive values.\ \pause This means that it has a minimum value $s > 0$.\\
\pause
Since $\vec x$ is an accumulation point, $B(\vec x; s)\cap A$ has an element $\vec y$ different from $\vec x$.\\
\pause
However, $\vec y\in C$ and $|\vec y - \vec x| < s$, contradicting the minimality of $s$.\\
\pause
We conclude that $C$ is infinite.
\end{proof}
\end{frame}


\begin{frame}{Closure of a set}
\begin{defn}
The \textbf{closure} of a set $A$ is the set $\overline{A}$ of all adherent points of $A$
\end{defn}
\pause
\begin{thm}[Apostol Theorem 3.18, 3.20, 3.22]
A set is closed if and only if it contains all of its adherent points (ie. $A= \overline{A}$), or equivalently if and only if $A$ contains all of its accumulation points. 
\end{thm}
\end{frame}

\begin{frame}{Closure of a set}
\begin{proof}
Suppose that $A$ is closed.\\
\pause
Then $\mathbb{R}^n\backslash A$ is open.\\
\pause
If $\vec x\in \mathbb{R}^n\backslash A$, then there exist $r > 0$ such that $B(\vec x; r)\subseteq \mathbb{R}^n\backslash A$.\\
\pause
This means that $B(\vec x; r)\cap A = \varnothing$.\\
\pause
Thus $\vec x$ is not an adherent point of $A$.\\
\pause
Thus every adherent point of $A$ is an element of $A$.\\
\pause
Conversely, suppse that $A$ contains all of its accumulation points.\\
\pause
If $\vec x\in \mathbb R^n\backslash A$, then $\vec x$ is not an accumulation point.\\
\pause
Therefore there exists $r > 0$ with $B(\vec x; r)\cap A = \varnothing$.\\
\pause
This implies that $B(\vec x; r)\subseteq \mathbb R^n\backslash A$.\\
\pause
Thus $\vec x$ is an interior point of $\mathbb R^n\backslash A$.\\
\pause
Since $\vec x$ was an arbitrary element of $\mathbb R^n\backslash A$, this proves $\mathbb R^n\backslash A$ is open.\\
\pause
Therefore $A$ is closed.
\end{proof}
\end{frame}

\begin{frame}{Closures are closed}
\begin{thm}
If $A\subseteq \mathbb{R}^n$ is any set, then $\overline{A}$ is closed.
\end{thm}
\begin{proof}
Exercise!
\end{proof}
\end{frame}

\begin{frame}{Closed balls are closed}
\begin{prob}
Show that the closed ball
$$\overline B(\vec x; r) = \{\vec y\in \mathbb{R}^n: |\vec x-\vec y| \leq r\}$$
is a closed set.
\end{prob}
\end{frame}


\end{document}


